\chapter{Movement Regularity}

\noindent\textbf{Movement Regularity} estimates \textit{how rhythmic and self-repeating
the body's motion is over a short window} --- a smoothness score, not an intensity or a
posture. It is theoretically distinct from \textbf{Activity State} (which measures
\textit{how much} motion and of \textit{what type}) and from \textbf{Postural State}
(which measures orientation relative to gravity). A person can be highly active yet
irregular (stumbling, fidgeting) or barely active yet perfectly regular (a slow steady
walk); regularity is the axis that separates these, so it is a genuinely non-redundant
construct rather than a re-scaling of activity intensity.

\begin{callout}
\textbf{Scope.} This score reports \textit{whether} motion repeats cleanly, not \textit{why}
it does or does not. A low score cannot by itself distinguish a mechanical cause (rough
terrain, an obstacle) from a physiological one (tremor, agitation, fatigue). The score is
only defined while the body is actually moving rhythmically; when the person is still, or
when too few motion cycles fall inside the window, the axis returns \textbf{null} rather
than a misleadingly ``perfect'' value. Resolving the cause of irregularity requires the
Activity State context and the downstream heart-rate / affective layer.
\end{callout}

\section{Theoretical Basis}

The score is the normalized \emph{unbiased autocorrelation} of the acceleration-magnitude
signal at its dominant period. This is the trunk-accelerometry gait-regularity index of
\textbf{Moe-Nilssen \& Helbostad (2004)}: for a walking signal, the value of the normalized
autocorrelation at the stride lag measures how closely one stride reproduces the next, and
the authors note it is cheap enough for portable devices. The parts of the formula
\emph{directly evidenced} by that work are the unbiased-autocorrelation estimator and the
use of its dominant-period peak height as the regularity index. The remaining choices are
\emph{engineering adaptations}, flagged as such below: the clamp to $[0,1]$, the movement
and cycle-count gates, the pocket-placement lag band, and the two-factor confidence.

\begin{itemize}[leftmargin=1.4em]
  \item \textbf{Acceleration magnitude $|a_\text{body}|$} --- gravity-removed body
  acceleration reduced to its Euclidean norm, making the feature orientation-invariant
  (the same rhythm reads identically regardless of how the phone is rotated in the pocket).
  \item \textbf{Unbiased autocorrelation} --- the $1/(N-k)$ normalization corrects the
  end-of-window sample deficit so that longer lags are not artificially damped
  (Moe-Nilssen \& Helbostad, 2004).
  \item \textbf{Smoothness alternatives (not used, cited for completeness)} --- spectral
  arc length SPARC (Balasubramanian et al., 2012) and log dimensionless jerk (LDLJ) are
  established smoothness measures; autocorrelation is chosen here because it is
  period-aware and matches the gait literature.
  \item \textbf{Affective relevance} --- reduced or irregular ambulatory movement is
  associated with agitation and anxiety in actigraphy studies, which motivates exposing
  regularity to the downstream affective read.
\end{itemize}

\section{Input Signals}

\begin{center}
\begin{tabular}{@{}l l >{\raggedright\arraybackslash}p{7.4cm}@{}}
\toprule
\textbf{Signal} & \textbf{Symbol} & \textbf{Description} \\
\midrule
Body acceleration      & $a_\text{body}$ & Gravity-removed 3-axis acceleration (m/s$^2$);
                                            linear-accel sensor, or raw accel low-passed at 0.3\,Hz \\
Acceleration magnitude & $s$             & $|a_\text{body}|/g$, orientation-invariant motion size (g) \\
Dominant lag           & $k^\star$       & Period (samples) of the strongest self-repeat in the window \\
Activity state         & \texttt{ACT}    & Context gate: regularity is only meaningful while walking/running \\
\bottomrule
\end{tabular}
\end{center}

\noindent Window $N = 250$ samples ($5$\,s at $F_s = 50$\,Hz), 50\% overlap. No personal
baseline is required: the score is self-referential within the window.

\section{Rule Formula}

\begin{rulebox}{RULE-MR}{Movement Regularity Score}
\textbf{Step 1 --- Magnitude and detrend:}
\[ s(t) = \frac{|a_\text{body}(t)|}{g}, \qquad x(t) = s(t) - \overline{s}, \qquad
   \text{movement} = \std(s). \]

\textbf{Step 2 --- Unbiased, normalized autocorrelation} (lags $k = 0\dots K$, $K=100$):
\[ r(k) = \frac{1}{N-k}\sum_{t=0}^{N-1-k} x(t)\,x(t+k), \qquad
   \rho(k) = \frac{r(k)}{r(0)}. \]

\textbf{Step 3 --- Dominant period} (argmax over the stride band, not the first peak):
\[ k^\star = \operatorname*{arg\,max}_{k\in[10,100]} \rho(k), \qquad
   n_\text{cyc} = \frac{N}{k^\star}. \]

\textbf{Step 4 --- Score:}
\[ R = \clamp\big(\rho(k^\star),\,0,\,1\big). \]

\textbf{Step 5 --- Validity gate} (return \textbf{null} instead of a score if):
\[ \text{movement} < 0.051\text{ g}\ \ (\text{low\_movement}) \quad\text{or}\quad
   n_\text{cyc} < 3.5\ \ (\text{insufficient\_cycles}). \]

\boxform{ R \in [0,1] \ \text{ if valid, else null} }

\smallskip
\textit{Example:} steady pocket walk, $\std(s)=0.37$\,g, $k^\star=55$ ($\approx0.9$\,Hz
stride), $\rho(k^\star)=0.71$, $n_\text{cyc}=250/55=4.5$. Both gates pass
$\Rightarrow R = 0.71$ (``rhythmic'').
\end{rulebox}

\section{Parameter Notes}

\noindent\textbf{Regularity estimator} {\footnotesize\textsc{[theory]}}\textbf{.}
The score itself --- the normalized unbiased autocorrelation at the dominant period --- is
taken directly from the Moe-Nilssen \& Helbostad (2004) trunk-gait index. It is the one number
here that comes straight from published theory; everything else is a gate or a confidence
around it.

\noindent\textbf{Lag band $[10,100]$ samples} {\footnotesize\textsc{[conv]}}\textbf{.}
At $50$\,Hz this covers stride periods $0.2$--$2.0$\,s ($0.5$--$5$\,Hz), the plausible human
cadence range. Taking the \emph{argmax} over the whole band (rather than the first
prominence peak) is deliberate: on pocket data the autocorrelation has strong sub-stride
harmonics, and the first peak grabs a half-stride and inverts the rank ordering. This is the
FIX~2 correction re-derived during the chest$\rightarrow$pocket convergence test.

\noindent\textbf{Movement floor $0.051$\,g} {\footnotesize\textsc{[design]}}\textbf{.}
Below this, magnitude variation is sensor noise, not motion, and the autocorrelation of
noise produces a spurious high peak. The value is the chest-derived $0.5$\,m/s$^2$ floor
re-expressed in g ($0.5/9.81$); expressing it in stream units was FIX~1 of the convergence
test.

\noindent\textbf{Minimum cycles $3.5$ strides} {\footnotesize\textsc{[conv]}}\textbf{.}
Fewer than $\sim$3--4 repeats in the window makes a ``regularity'' estimate statistically
meaningless. $3.5$ stride-periods is the gait-literature adequacy minimum, re-expressed for
the argmax-selects-stride convention on pocket placement.

\noindent\textbf{Prominence saturation $P_\text{full}=1.5$} {\footnotesize\textsc{[design]}}\textbf{.}
The peak prominence at which the confidence prominence term reaches $1$; a chosen saturation
point, not a measured quantity.

\noindent\textbf{Fusion method.}
There is \emph{no} multi-signal fusion in the score itself --- it is a single autocorrelation
read. Fusion appears only in the \emph{confidence} (\S\,Confidence), which combines two
\emph{independent} adequacy factors (enough cycles, and a prominent peak) as a geometric
mean, so that either factor being weak pulls confidence down. The score's \emph{separation}
(not an absolute accuracy) was validated on PAMAP2/MotionSense; the gates and confidence are
\textsc{design} choices tuned to keep noise from scoring as rhythm. Tag meanings are defined
in the Key Terms front-matter.

\section{Confidence}

\noindent\textbf{Step C1 --- Cycle adequacy} (how far past the minimum the window sits):
\[ c_\text{cyc} = \clamp\!\left(\frac{n_\text{cyc}-3.5}{5.0-3.5},\,0,\,1\right). \]

\noindent\textbf{Step C2 --- Peak prominence} ($p$ = prominence of the peak at $k^\star$;
$1.5$ = saturation):
\[ c_\text{prom} = \clamp\!\left(\frac{p}{1.5},\,0,\,1\right). \]

\noindent\textbf{Step C3 --- Final confidence} (geometric mean of the two):
\boxform{ \text{conf} = \sqrt{\,c_\text{cyc}\cdot c_\text{prom}\,} }

\noindent\emph{Confidence ceiling / caveat.} This axis reports no labelled accuracy by
design --- it is validated by \emph{separation} (rank-ordering rhythmic above irregular
above still), not against a ground-truth ``regularity'' label, which does not exist in the
datasets. Confidence therefore expresses \emph{internal signal quality} (enough clean
cycles), not calibrated correctness, and should not be read as a probability.

\section{Interpretation Scale}

\begin{center}
\begin{tabular}{@{}lll@{}}
\toprule
\textbf{Score} & \textbf{Level} & \textbf{Meaning} \\
\midrule
0.0--0.2 & Erratic          & Motion present but not self-repeating (stumbling, fidgeting) \\
0.2--0.4 & Low regularity   & Weak, inconsistent rhythm \\
0.4--0.6 & Moderate         & Discernible but imperfect rhythm \\
0.6--0.8 & Rhythmic         & Clear, steady periodic motion (typical walking) \\
0.8--1.0 & Highly regular   & Metronomic, strongly self-repeating (steady running/pace) \\
\midrule
\textit{null} & Not applicable & Still, or too few cycles --- no rhythm to score \\
\bottomrule
\end{tabular}
\end{center}

\section{Context Applicability}

\begin{callout}
\textbf{Target context:} on-body (pocket / trunk) accelerometry during \emph{sustained
locomotion}. The score is most reliable when the Activity State is walking or running.
\end{callout}

\begin{center}
\begin{tabular}{@{}>{\raggedright\arraybackslash}p{5.6cm} c c >{\raggedright\arraybackslash}p{3.6cm}@{}}
\toprule
\textbf{Context} & \textbf{$|a_\text{body}|$} & \textbf{cycles $\geq 3.5$} & \textbf{Reliable?} \\
\midrule
Walking / running (pocket or trunk) & Yes & Yes & Yes --- full \\
Cycling / non-gait rhythmic motion  & Yes & Partial & Partial (period may fall outside band) \\
Standing / sitting / lying          & No  & No  & No --- returns null \\
\bottomrule
\end{tabular}
\end{center}

\noindent\textbf{Scientific boundary.} Even with ideal data, the score cannot attribute
\emph{cause}: it measures the presence of rhythm, not its origin, so it approximates
``movement smoothness'' rather than directly measuring any affective or neuromotor state.
On PAMAP2 (chest) the score separates rhythmic from irregular-but-moving activity, so it is
not merely a movement detector; on the pocket, only the still-versus-rhythmic contrast is
confirmed so far, and rhythmic-versus-irregular-moving on the pocket is owed to own-phone data.

\section{Worked Example: Steady Commute Walk}

\noindent\textit{Scenario.} A user walks to a bus stop at a relaxed, even pace, phone in the
front trouser pocket. Over one $5$\,s window the pipeline extracts:

\begin{center}
\begin{tabular}{@{}ll@{}}
\toprule
\textbf{Signal} & \textbf{Value} \\
\midrule
$\std(s)$ (movement)        & $0.37$\,g \\
Dominant lag $k^\star$      & $55$ samples ($\approx 0.9$\,Hz stride) \\
$\rho(k^\star)$             & $0.71$ \\
Peak prominence $p$         & $0.60$ \\
\bottomrule
\end{tabular}
\end{center}

\noindent\textbf{Step-by-step:}
\begin{align*}
n_\text{cyc} &= 250/55 = 4.5 \;(\geq 3.5,\ \text{movement } 0.37>0.051 \Rightarrow \text{valid})\\
R            &= \clamp(0.71,0,1) = \mathbf{0.71}\\
c_\text{cyc} &= \clamp\big((4.5-3.5)/1.5,0,1\big) = 0.67\\
c_\text{prom}&= \clamp(0.60/1.5,0,1) = 0.40\\
\text{conf}  &= \sqrt{0.67\cdot 0.40} = \mathbf{0.52}
\end{align*}

\begin{callout}
\textbf{Interpretation:} $R = 0.71$ falls in the \emph{Rhythmic} band --- a clean, steady
walking cadence. The dominant driver is the high autocorrelation peak ($\rho=0.71$);
confidence is moderate ($0.52$), held down by the modest peak prominence rather than by
too few cycles.
\end{callout}

\section{Implementation Notes}

\begin{itemize}[leftmargin=1.4em]
  \item \textbf{Window.} $250$ samples ($5$\,s) at $50$\,Hz, 50\% overlap; magnitude is
  read from the same rolling buffer the other 5\,s axes use.
  \item \textbf{Fallback behavior.} No accel/linear stream $\rightarrow$ null. A window that
  is $>50\%$ bad samples ($\mathrm{NaN}$ or $>16$\,g), or whose first sample is bad,
  $\rightarrow$ null (\texttt{window\_unreliable}); shorter gaps are forward-filled. If
  \texttt{scipy} is unavailable the score still computes; only the prominence confidence
  factor degrades.
  \item \textbf{Relationship to other chapters.} Gated by \textbf{Activity State} --- a
  regularity number is only surfaced while \texttt{ACT} indicates locomotion; on its own,
  regularity cannot tell moving-irregular from still, which is why the two axes are read
  together.
  \item \textbf{Validation status.} No labelled accuracy by design (there is no ground-truth
  regularity label), so validation is by \emph{separation}. On PAMAP2 chest, rhythmic vs
  irregular ROC-AUC $0.99$ (per-subject median $0.99$, all 8 evaluable subjects $\geq 0.97$).
  It transfers \emph{frozen} to two further datasets and sensor types (table below).
  \textbf{And it is regularity, not movement:} restricted to \emph{moving} windows --- so
  ``moving vs still'' cannot explain the split --- rhythmic vs irregular-but-moving
  (vacuuming, ironing) still separates at ROC-AUC $0.99$: vacuuming moves yet scores
  ${\sim}0.25$ while walking scores ${\sim}0.94$ (regularity correlates with movement only
  ${\sim}{+}0.5$). This movement-control exists only on PAMAP2; the pocket version is owed to
  own-phone data.
\end{itemize}

\begin{center}
\begin{tabular}{@{}>{\raggedright\arraybackslash}p{2.4cm} >{\raggedright\arraybackslash}p{3.6cm} >{\raggedright\arraybackslash}p{3.4cm} l@{}}
\toprule
\textbf{Dataset} & \textbf{Placement / sensor} & \textbf{Contrast} & \textbf{ROC-AUC} \\
\midrule
PAMAP2      & chest, research IMU        & rhythmic vs irregular & $0.99$ \\
MotionSense & pocket, phone              & rhythmic vs still     & $0.96$ \\
HHAR        & waist/hand, phone + watch  & walking vs still      & $0.98$ \\
\bottomrule
\end{tabular}
\end{center}

\noindent{\footnotesize MotionSense/HHAR have no irregular-but-moving activity, so
cross-dataset confirms rhythmic $>$ still only. HHAR must be resampled \emph{within gap-free
spans}: its heterogeneous sampling, interpolated across gaps, otherwise fakes periodicity.}

% ---- Formula flow diagram ----
\begin{figure}[h]
\centering
\begin{tikzpicture}[
  node distance=10mm and 14mm,
  box/.style={draw, dashed, rounded corners=1pt, align=center, inner sep=5pt, font=\small},
  op/.style={draw, rounded corners=1pt, align=center, inner sep=5pt, font=\small},
  resultbox/.style={draw, double, rounded corners=3pt, align=center, inner sep=6pt, font=\small\bfseries},
  gate/.style={draw, dashed, align=center, inner sep=4pt, font=\footnotesize},
  ar/.style={-{Latex[length=2mm]}}]

  \node[box] (in) {Body acceleration\\ $|a_\text{body}|/g$};
  \node[op, below=of in] (ac) {Unbiased autocorrelation\\ $\rho(k)$ \; (Steps 1--2)};
  \node[op, below=of ac] (peak) {Dominant period $k^\star$\\ $=\arg\max_{[10,100]}\rho$ \; (Step 3)};
  \node[gate, right=10mm of peak] (g) {Gates:\\ movement $\geq 0.051$g\\ $n_\text{cyc}\geq 3.5$};
  \node[resultbox, below=of peak] (score) {$R=\clamp(\rho(k^\star),0,1) \in [0,1]$};

  \draw[ar] (in) -- (ac);
  \draw[ar] (ac) -- (peak);
  \draw[ar] (peak) -- (score);
  \draw[ar] (g) -- (score) node[midway,above,font=\footnotesize] {else null};
\end{tikzpicture}
\caption*{\small Formula flow --- Movement Regularity: one orientation-invariant input,
period-aware autocorrelation, a validity gate, and a $[0,1]$ score.}
\end{figure}
